\documentclass[11pt,a4paper]{article}
\usepackage[T1]{fontenc}
\usepackage{lmodern,array,amsmath,amssymb,booktabs,enumitem,fancyhdr,lastpage,needspace,microtype}
\usepackage[margin=19mm,headheight=14pt,footskip=11mm]{geometry}
\usepackage[hidelinks]{hyperref}
\setlength{\parindent}{0pt}
\setlength{\parskip}{5pt}
\setlength{\emergencystretch}{2em}
\pagestyle{fancy}\fancyhf{}
\fancyhead[L]{\small CS182 --- Lecture 07}
\fancyhead[R]{\small Personal answer copy}
\fancyfoot[L]{\small Independent practice quiz}
\fancyfoot[R]{\small Page \thepage\ of \pageref{LastPage}}
\newcommand{\SetAnswer}[2]{\expandafter\gdef\csname QuizAnswer#1\endcsname{#2}}
\newcommand{\GetAnswer}[1]{\csname QuizAnswer#1\endcsname}
\newlength{\BlankAnswerHeight}\setlength{\BlankAnswerHeight}{16mm}
\newcommand{\AnswerBlank}{\leavevmode\par\vspace{\BlankAnswerHeight}}
% 本文件为独立可填写答题卷；用 pdfLaTeX 编译两次。
% 在 BEGIN TF ANSWERS 和 BEGIN MC ANSWERS 中填写短答案。
% 大题请搜索 BEGIN ANSWER Q25(a) 等注释；每个小问都有独立区域。
% 写完小问后删除该处的 \AnswerBlank，或将 BlankAnswerHeight 改为 0mm。
% 所有打印内容及你填写的回答请使用英文。中文只出现在源文件注释中。
% build_per_lecture.py 不会覆盖已经存在的 Quiz.tex，以保留你的作答。
% BEGIN STUDENT INFO / 个人信息
\newcommand{\StudentName}{}
\newcommand{\StudentID}{}
% END STUDENT INFO
% BEGIN TF ANSWERS / 判断题答题区（第二组花括号填 T 或 F）
\SetAnswer{TF01}{} % Question 1
\SetAnswer{TF02}{} % Question 2
\SetAnswer{TF03}{} % Question 3
\SetAnswer{TF04}{} % Question 4
\SetAnswer{TF05}{} % Question 5
\SetAnswer{TF06}{} % Question 6
\SetAnswer{TF07}{} % Question 7
\SetAnswer{TF08}{} % Question 8
\SetAnswer{TF09}{} % Question 9
\SetAnswer{TF10}{} % Question 10
\SetAnswer{TF11}{} % Question 11
\SetAnswer{TF12}{} % Question 12
% END TF ANSWERS
% BEGIN MC ANSWERS / 不定项选择答题区（填选项字母，如 BD；仅格式示例）
\SetAnswer{MC13}{} % Question 13
\SetAnswer{MC14}{} % Question 14
\SetAnswer{MC15}{} % Question 15
\SetAnswer{MC16}{} % Question 16
\SetAnswer{MC17}{} % Question 17
\SetAnswer{MC18}{} % Question 18
\SetAnswer{MC19}{} % Question 19
\SetAnswer{MC20}{} % Question 20
\SetAnswer{MC21}{} % Question 21
\SetAnswer{MC22}{} % Question 22
\SetAnswer{MC23}{} % Question 23
\SetAnswer{MC24}{} % Question 24
% END MC ANSWERS

\begin{document}
\begin{center}{\Large\bfseries CS182: Introduction to Machine Learning}\par{\Large\bfseries Lecture 07: Linear Discrimination}\par{\large Lecture Quiz and Answer Copy}\end{center}
\textbf{Scope:} Linear\_Discrimination.pdf. Independent practice material based on the official 2026 Fall lecture.\par\textbf{Total:} 100 marks. \quad \textbf{Suggested time:} 120 minutes.
\par\textbf{Name:} \underline{\makebox[60mm][l]{\StudentName}}\hfill\textbf{Student ID:} \underline{\makebox[45mm][l]{\StudentID}}
\begin{center}\begin{tabular}{lrrr}\toprule
Section & Questions & Each & Marks\\\midrule
A. Misconceptions: True/False & 12 & 1 & 12\\
B. Multiple response & 12 & 2 & 24\\
C. Proofs and formula derivations & 4 & 10 & 40\\
D. Applications & 2 & 12 & 24\\\midrule
Total & 30 & & 100\\\bottomrule\end{tabular}\end{center}
\textbf{Scoring.} In Section B each item has exactly four options, of which one,
two, or three are correct. The complete correct set earns 2 marks; a nonempty
proper subset earns 1 mark; any incorrect option or no selection earns 0.
T/F answers earn 1 mark only when correct. Show reasoning in Sections C and D;
part marks are available for valid intermediate work. Equivalent correct methods receive credit.
\par\medskip\textbf{Instructions.} Answer in English. Use natural logarithms.
Exact values and justified approximations are accepted. All vectors are columns
unless stated otherwise. Use local assumptions, label conventions, and tie rules.
\par\medskip\textbf{Answer sheet --- Section A}
\begin{center}\renewcommand{\arraystretch}{1.15}\begin{tabular}{|c|c|c|c|c|c|}
\hline
1 & 2 & 3 & 4 & 5 & 6\\\hline
\makebox[20mm]{\rule{0pt}{5mm}\GetAnswer{TF01}} & \makebox[20mm]{\rule{0pt}{5mm}\GetAnswer{TF02}} & \makebox[20mm]{\rule{0pt}{5mm}\GetAnswer{TF03}} & \makebox[20mm]{\rule{0pt}{5mm}\GetAnswer{TF04}} & \makebox[20mm]{\rule{0pt}{5mm}\GetAnswer{TF05}} & \makebox[20mm]{\rule{0pt}{5mm}\GetAnswer{TF06}}\\\hline
\end{tabular}\end{center}
\begin{center}\renewcommand{\arraystretch}{1.15}\begin{tabular}{|c|c|c|c|c|c|}
\hline
7 & 8 & 9 & 10 & 11 & 12\\\hline
\makebox[20mm]{\rule{0pt}{5mm}\GetAnswer{TF07}} & \makebox[20mm]{\rule{0pt}{5mm}\GetAnswer{TF08}} & \makebox[20mm]{\rule{0pt}{5mm}\GetAnswer{TF09}} & \makebox[20mm]{\rule{0pt}{5mm}\GetAnswer{TF10}} & \makebox[20mm]{\rule{0pt}{5mm}\GetAnswer{TF11}} & \makebox[20mm]{\rule{0pt}{5mm}\GetAnswer{TF12}}\\\hline
\end{tabular}\end{center}
\par\medskip\textbf{Answer sheet --- Section B}
\begin{center}\renewcommand{\arraystretch}{1.15}\begin{tabular}{|c|c|c|c|c|c|}
\hline
13 & 14 & 15 & 16 & 17 & 18\\\hline
\makebox[20mm]{\rule{0pt}{5mm}\GetAnswer{MC13}} & \makebox[20mm]{\rule{0pt}{5mm}\GetAnswer{MC14}} & \makebox[20mm]{\rule{0pt}{5mm}\GetAnswer{MC15}} & \makebox[20mm]{\rule{0pt}{5mm}\GetAnswer{MC16}} & \makebox[20mm]{\rule{0pt}{5mm}\GetAnswer{MC17}} & \makebox[20mm]{\rule{0pt}{5mm}\GetAnswer{MC18}}\\\hline
\end{tabular}\end{center}
\begin{center}\renewcommand{\arraystretch}{1.15}\begin{tabular}{|c|c|c|c|c|c|}
\hline
19 & 20 & 21 & 22 & 23 & 24\\\hline
\makebox[20mm]{\rule{0pt}{5mm}\GetAnswer{MC19}} & \makebox[20mm]{\rule{0pt}{5mm}\GetAnswer{MC20}} & \makebox[20mm]{\rule{0pt}{5mm}\GetAnswer{MC21}} & \makebox[20mm]{\rule{0pt}{5mm}\GetAnswer{MC22}} & \makebox[20mm]{\rule{0pt}{5mm}\GetAnswer{MC23}} & \makebox[20mm]{\rule{0pt}{5mm}\GetAnswer{MC24}}\\\hline
\end{tabular}\end{center}
\newpage\section*{Section A. True/False \hfill 12 marks}
\par\noindent\begin{minipage}{\linewidth}
\textbf{1.} A multiclass argmax classifier with affine scores can have a nonconvex class region even when there are no ties.
\end{minipage}\par\vspace{3mm}
\par\noindent\begin{minipage}{\linewidth}
\textbf{2.} For any fixed \(w\) and \(w_0\), the set \(\{x:w^{\mathsf T}x+w_0=0\}\) is a separating hyperplane of dimension \(d-1\).
\end{minipage}\par\vspace{3mm}
\par\noindent\begin{minipage}{\linewidth}
\textbf{3.} Multiplying both the weights and intercept of a binary linear score by a positive constant preserves the predicted labels but generally changes its sigmoid probabilities.
\end{minipage}\par\vspace{3mm}
\par\noindent\begin{minipage}{\linewidth}
\textbf{4.} The batch algorithm shown in the lecture accumulates errors over all samples before updating the weight vector.
\end{minipage}\par\vspace{3mm}
\par\noindent\begin{minipage}{\linewidth}
\textbf{5.} For finite binary logits and no numerical rounding, sigmoid produces probabilities strictly between zero and one.
\end{minipage}\par\vspace{3mm}
\par\noindent\begin{minipage}{\linewidth}
\textbf{6.} For one-hot softmax cross-entropy, a sample generally contributes a gradient to the weights of classes other than its true class.
\end{minipage}\par\vspace{3mm}
\par\noindent\begin{minipage}{\linewidth}
\textbf{7.} Training \(K\) independent one-vs.-rest sigmoid models automatically makes their \(K\) outputs sum to one.
\end{minipage}\par\vspace{3mm}
\par\noindent\begin{minipage}{\linewidth}
\textbf{8.} Adding the same input-dependent affine function to every softmax class score leaves every predicted class probability unchanged.
\end{minipage}\par\vspace{3mm}
\par\noindent\begin{minipage}{\linewidth}
\textbf{9.} Two learning procedures that both produce a linear class boundary must estimate their parameters by the same statistical objective.
\end{minipage}\par\vspace{3mm}
\par\noindent\begin{minipage}{\linewidth}
\textbf{10.} When training a one-vs.-one classifier for classes \(C_i\) and \(C_j\), examples from other classes are excluded from that particular binary training problem.
\end{minipage}\par\vspace{3mm}
\par\noindent\begin{minipage}{\linewidth}
\textbf{11.} A standard linear logistic-regression model requires every input feature to have a Gaussian class-conditional distribution.
\end{minipage}\par\vspace{3mm}
\par\noindent\begin{minipage}{\linewidth}
\textbf{12.} Convexity of unregularized logistic cross-entropy guarantees a finite maximum-likelihood parameter vector on every strictly linearly separable dataset.
\end{minipage}\par\vspace{3mm}
\newpage\section*{Section B. Multiple Response \hfill 24 marks}
\textit{Select all correct options. Each question has 1--3 correct options.}\par
\par\noindent\begin{minipage}{\linewidth}
% BEGIN QUESTION Q13 / 题目与选项区
\textbf{13.} Let a batch contain four samples. Compare summed cross-entropy and average cross-entropy at exactly the same weights. Which statements are correct?
\begin{enumerate}[label=(\Alph*),leftmargin=8mm,itemsep=2pt,topsep=2pt]
\item At the same learning rate, the average-loss gradient step is one quarter of the summed-loss step.
\item Choosing the average-loss learning rate four times as large reproduces the summed-loss update.
\item Zero initialization forces the binary logistic gradient to be zero for every labeled dataset.
\item Updating weights after every sample is necessarily identical to accumulating the four gradients at fixed weights.
\end{enumerate}
% END QUESTION Q13

\end{minipage}\par\vspace{3mm}
\par\noindent\begin{minipage}{\linewidth}
% BEGIN QUESTION Q14 / 题目与选项区
\textbf{14.} Three independently trained affine scores at one test point are \((-2,-1,-3)\). The specified decision rule is argmax over these scores, with no rejection option. Which statements are correct?
\begin{enumerate}[label=(\Alph*),leftmargin=8mm,itemsep=2pt,topsep=2pt]
\item Adding ten to every score changes the argmax class.
\item The point must be rejected because every score is negative.
\item No label can be assigned until the scores have been individually transformed to probabilities.
\item The predicted class is \(C_2\).
\end{enumerate}
% END QUESTION Q14

\end{minipage}\par\vspace{3mm}
\par\noindent\begin{minipage}{\linewidth}
% BEGIN QUESTION Q15 / 题目与选项区
\textbf{15.} Correct decisions cost zero. Predicting \(C_1\) when the truth is \(C_2\) costs three; the reverse error costs two. At a point with \(y=P(C_1\mid x)=0.55\), which statements are correct?
\begin{enumerate}[label=(\Alph*),leftmargin=8mm,itemsep=2pt,topsep=2pt]
\item The risk of predicting \(C_2\) at the point is \(0.9\).
\item The rule must predict \(C_1\) at the given point because \(0.55>0.5\).
\item The risk-minimizing rule predicts \(C_1\) when \(y>3/5\).
\item The equivalent logit threshold is \(\log(2/3)\).
\end{enumerate}
% END QUESTION Q15

\end{minipage}\par\vspace{3mm}
\par\noindent\begin{minipage}{\linewidth}
% BEGIN QUESTION Q16 / 题目与选项区
\textbf{16.} Consider adding \(\lambda\|\theta\|_2^2/2\), with \(\lambda>0\), to a finite-sample linear binary logistic loss, where \(\theta\) includes the intercept. Which statements are correct?
\begin{enumerate}[label=(\Alph*),leftmargin=8mm,itemsep=2pt,topsep=2pt]
\item The added Hessian is \(\lambda I\).
\item The penalty weight can be selected using validation data.
\item The penalty necessarily sets some fitted coefficients exactly to zero.
\item The penalized objective has a unique finite minimizer.
\end{enumerate}
% END QUESTION Q16

\end{minipage}\par\vspace{3mm}
\newpage\section*{Section B. Multiple Response \hfill 24 marks}
\textit{Select all correct options. Each question has 1--3 correct options.}\par
\par\noindent\begin{minipage}{\linewidth}
% BEGIN QUESTION Q17 / 题目与选项区
\textbf{17.} For \(\sigma(a)=1/(1+e^{-a})\), which statements are correct?
\begin{enumerate}[label=(\Alph*),leftmargin=8mm,itemsep=2pt,topsep=2pt]
\item The derivative equals zero at every finite positive logit.
\item The largest derivative value is \(1/4\), attained at \(a=0\).
\item \(\sigma(a+b)=\sigma(a)+\sigma(b)\) for all real \(a,b\).
\item \(\sigma(\log9)=0.9\).
\end{enumerate}
% END QUESTION Q17

\end{minipage}\par\vspace{3mm}
\par\noindent\begin{minipage}{\linewidth}
% BEGIN QUESTION Q18 / 题目与选项区
\textbf{18.} A binary logistic model predicts \(y=1/4\) for a sample with label \(r=1\) and feature vector \(x=(1,-2)^{\mathsf T}\). Which statements are correct?
\begin{enumerate}[label=(\Alph*),leftmargin=8mm,itemsep=2pt,topsep=2pt]
\item The intercept gradient is \(-3/4\).
\item The negative-gradient update decreases the intercept when the learning rate is positive.
\item The feature-weight gradient is \((-3/4,-3/2)^{\mathsf T}\).
\item The sample cross-entropy is \(\log4\).
\end{enumerate}
% END QUESTION Q18

\end{minipage}\par\vspace{3mm}
\par\noindent\begin{minipage}{\linewidth}
% BEGIN QUESTION Q19 / 题目与选项区
\textbf{19.} Two Gaussian classes have means \(1\) and \(-1\), shared variance two, and equal priors. Which statements are correct?
\begin{enumerate}[label=(\Alph*),leftmargin=8mm,itemsep=2pt,topsep=2pt]
\item This derivation proves that all logistic input distributions must be Gaussian.
\item The posterior probability of \(C_1\) is \(\sigma(x)\).
\item Changing only the priors, with means and variance fixed, changes only the logit intercept.
\item Two distinct Gaussian variances always yield the same linear logit.
\end{enumerate}
% END QUESTION Q19

\end{minipage}\par\vspace{3mm}
\par\noindent\begin{minipage}{\linewidth}
% BEGIN QUESTION Q20 / 题目与选项区
\textbf{20.} At a point where softmax probabilities are \((0.2,0.5,0.3)\), which statements are correct?
\begin{enumerate}[label=(\Alph*),leftmargin=8mm,itemsep=2pt,topsep=2pt]
\item \(\partial y_2/\partial a_1=-0.1\).
\item Every column of the Jacobian sums to one.
\item \(\partial y_2/\partial a_2=0.5\).
\item The Jacobian is diagonal because each numerator uses one score.
\end{enumerate}
% END QUESTION Q20

\end{minipage}\par\vspace{3mm}
\newpage\section*{Section B. Multiple Response \hfill 24 marks}
\textit{Select all correct options. Each question has 1--3 correct options.}\par
\par\noindent\begin{minipage}{\linewidth}
% BEGIN QUESTION Q21 / 题目与选项区
\textbf{21.} Consider \(g(x)=x_1^2+x_2^2-1\), with \(C_1\) predicted for \(g(x)>0\). Which statements are correct?
\begin{enumerate}[label=(\Alph*),leftmargin=8mm,itemsep=2pt,topsep=2pt]
\item The score is affine in the transformed feature \(z=x_1^2+x_2^2\).
\item The boundary is a straight line in the original input plane because the coefficients enter linearly.
\item The boundary is a unit circle in the original input plane.
\item The model requires \(x\) to have a Gaussian distribution.
\end{enumerate}
% END QUESTION Q21

\end{minipage}\par\vspace{3mm}
\par\noindent\begin{minipage}{\linewidth}
% BEGIN QUESTION Q22 / 题目与选项区
\textbf{22.} Let \(w=(4,-3)^{\mathsf T}\), \(w_0=-5\), and \(x=(2,0)^{\mathsf T}\). Which statements are correct?
\begin{enumerate}[label=(\Alph*),leftmargin=8mm,itemsep=2pt,topsep=2pt]
\item The ordinary distance from the origin to the boundary equals minus one.
\item The Euclidean distance from \(x\) to the boundary equals three.
\item The score at \(x\) equals three.
\item The orthogonal projection is \((38/25,9/25)^{\mathsf T}\).
\end{enumerate}
% END QUESTION Q22

\end{minipage}\par\vspace{3mm}
\par\noindent\begin{minipage}{\linewidth}
% BEGIN QUESTION Q23 / 题目与选项区
\textbf{23.} Three softmax logits are \((0,\log3,\log2)\). Which statements are correct?
\begin{enumerate}[label=(\Alph*),leftmargin=8mm,itemsep=2pt,topsep=2pt]
\item The second class has probability \(1/2\).
\item Adding seven to every logit preserves all three probabilities.
\item The predicted class must be rejected because its probability does not strictly exceed \(1/2\).
\item The posterior odds of \(C_2\) against \(C_3\) equal \(3/2\).
\end{enumerate}
% END QUESTION Q23

\end{minipage}\par\vspace{3mm}
\par\noindent\begin{minipage}{\linewidth}
% BEGIN QUESTION Q24 / 题目与选项区
\textbf{24.} A four-class softmax model has two original input features and an intercept per class. Which statements are correct?
\begin{enumerate}[label=(\Alph*),leftmargin=8mm,itemsep=2pt,topsep=2pt]
\item Adding the same three-component augmented weight vector to every class changes the probabilities.
\item Fixing the reference-class score to zero leaves nine free coefficients.
\item The symmetric parameterization has twelve coefficients.
\item For a row-weight matrix \(B\in\mathbb R^{4\times3}\), the gradient has the same \(4\times3\) shape.
\end{enumerate}
% END QUESTION Q24

\end{minipage}\par\vspace{3mm}
\newpage\section*{Section C. Proofs and Formula Derivations}
\textbf{25. Derive the geometry rather than treating the score as a distance}\hfill\textbf{[10 marks]}\par\medskip
\begin{enumerate}[label=(\alph*),leftmargin=8mm,itemsep=3pt]
\item Let \(g(x)=w^{\mathsf T}x+b\), with \(w\ne0\). Prove that \(w\) is normal to its zero hyperplane by comparing two points on that hyperplane.\hfill\textbf{[2 marks]}
\item Write a point as \(x=x_p+r w/\|w\|_2\), where \(x_p\) is its orthogonal projection. Derive both the signed-distance formula and an explicit formula for \(x_p\). Explain the corresponding unsigned distance.\hfill\textbf{[4 marks]}
\item Prove what happens to the hyperplane, signed distance, and predicted class when \((w,b)\) is multiplied by \(c>0\). Does the corresponding sigmoid probability remain fixed? Include the degenerate case \(w=0\).\hfill\textbf{[4 marks]}
\end{enumerate}\par\medskip\hrule\medskip\textbf{Student response}\par
% BEGIN ANSWER Q25 / 第25题答题区
\begin{enumerate}[label=(\alph*),leftmargin=8mm,itemsep=3pt]
\item
% >>> BEGIN ANSWER Q25(a)
% 在这里写英文推导或证明；完成后可删除下方空白占位命令。
\AnswerBlank
% <<< END ANSWER Q25(a)
\item
% >>> BEGIN ANSWER Q25(b)
% 在这里写英文推导或证明；完成后可删除下方空白占位命令。
\AnswerBlank
% <<< END ANSWER Q25(b)
\item
% >>> BEGIN ANSWER Q25(c)
% 在这里写英文推导或证明；完成后可删除下方空白占位命令。
\AnswerBlank
% <<< END ANSWER Q25(c)
\end{enumerate}
% END ANSWER Q25
\newpage\section*{Section C. Proofs and Formula Derivations}
\textbf{26. Derive the LDA-to-logistic bridge with its assumptions}\hfill\textbf{[10 marks]}\par\medskip
\begin{enumerate}[label=(\alph*),leftmargin=8mm,itemsep=3pt]
\item Assume \(x\mid C_i\sim\mathcal N(\mu_i,\Sigma)\) for two classes, with shared positive-definite \(\Sigma\) and positive priors \(\pi_i\). Derive \(\log[P(C_1\mid x)/P(C_2\mid x)]=w^{\mathsf T}x+b\), explicitly finding \(w,b\).\hfill\textbf{[5 marks]}
\item Solve \(a=\log[y/(1-y)]\) for \(y\), and prove the equivalence between the equal-cost posterior threshold and the zero-score threshold.\hfill\textbf{[2 marks]}
\item Identify the canceled term that generally remains with unequal class covariances. Explain why the Gaussian derivation does not impose a Gaussian-input requirement on discriminatively fitted logistic regression.\hfill\textbf{[3 marks]}
\end{enumerate}\par\medskip\hrule\medskip\textbf{Student response}\par
% BEGIN ANSWER Q26 / 第26题答题区
\begin{enumerate}[label=(\alph*),leftmargin=8mm,itemsep=3pt]
\item
% >>> BEGIN ANSWER Q26(a)
% 在这里写英文推导或证明；完成后可删除下方空白占位命令。
\AnswerBlank
% <<< END ANSWER Q26(a)
\item
% >>> BEGIN ANSWER Q26(b)
% 在这里写英文推导或证明；完成后可删除下方空白占位命令。
\AnswerBlank
% <<< END ANSWER Q26(b)
\item
% >>> BEGIN ANSWER Q26(c)
% 在这里写英文推导或证明；完成后可删除下方空白占位命令。
\AnswerBlank
% <<< END ANSWER Q26(c)
\end{enumerate}
% END ANSWER Q26
\newpage\section*{Section C. Proofs and Formula Derivations}
\textbf{27. Derive Bernoulli cross-entropy, its gradient, and its limits}\hfill\textbf{[10 marks]}\par\medskip
\begin{enumerate}[label=(\alph*),leftmargin=8mm,itemsep=3pt]
\item For conditionally independent labels \(r^t\in\{0,1\}\) with \(y^t=\sigma(\theta^{\mathsf T}\widetilde x^t)\), derive the conditional likelihood and its negative log-likelihood. Here \(\widetilde x^t\) contains the intercept feature.\hfill\textbf{[3 marks]}
\item Use the chain rule to derive \(\nabla_\theta E=\sum_t(y^t-r^t)\widetilde x^t\). Then derive the Hessian and prove convexity without treating sigmoid itself as a linear function.\hfill\textbf{[5 marks]}
\item For the two observations \((x,r)=(-2,0),(2,1)\), set intercept zero and slope \(c>0\). Prove that the unregularized minimum need not be attained at a finite parameter, although the classification error is already zero.\hfill\textbf{[2 marks]}
\end{enumerate}\par\medskip\hrule\medskip\textbf{Student response}\par
% BEGIN ANSWER Q27 / 第27题答题区
\begin{enumerate}[label=(\alph*),leftmargin=8mm,itemsep=3pt]
\item
% >>> BEGIN ANSWER Q27(a)
% 在这里写英文推导或证明；完成后可删除下方空白占位命令。
\AnswerBlank
% <<< END ANSWER Q27(a)
\item
% >>> BEGIN ANSWER Q27(b)
% 在这里写英文推导或证明；完成后可删除下方空白占位命令。
\AnswerBlank
% <<< END ANSWER Q27(b)
\item
% >>> BEGIN ANSWER Q27(c)
% 在这里写英文推导或证明；完成后可删除下方空白占位命令。
\AnswerBlank
% <<< END ANSWER Q27(c)
\end{enumerate}
% END ANSWER Q27
\newpage\section*{Section C. Proofs and Formula Derivations}
\textbf{28. Derive softmax normalization and the coupled learning rule}\hfill\textbf{[10 marks]}\par\medskip
\begin{enumerate}[label=(\alph*),leftmargin=8mm,itemsep=3pt]
\item Take \(C_K\) as reference and assume \(\log[P(C_i\mid x)/P(C_K\mid x)]=a_i\) for \(i<K\). Derive all class probabilities from normalization, then express them in symmetric softmax form.\hfill\textbf{[3 marks]}
\item Derive \(\partial y_i/\partial a_j\), then derive \(\partial\ell/\partial a_j=y_j-r_j\) for one-hot cross-entropy. Give feature-weight and intercept gradients for affine scores.\hfill\textbf{[4 marks]}
\item Prove softmax invariance to a common affine score shift, give the free parameter count after fixing a reference class, and explain why a nontrue class still receives a gradient.\hfill\textbf{[3 marks]}
\end{enumerate}\par\medskip\hrule\medskip\textbf{Student response}\par
% BEGIN ANSWER Q28 / 第28题答题区
\begin{enumerate}[label=(\alph*),leftmargin=8mm,itemsep=3pt]
\item
% >>> BEGIN ANSWER Q28(a)
% 在这里写英文推导或证明；完成后可删除下方空白占位命令。
\AnswerBlank
% <<< END ANSWER Q28(a)
\item
% >>> BEGIN ANSWER Q28(b)
% 在这里写英文推导或证明；完成后可删除下方空白占位命令。
\AnswerBlank
% <<< END ANSWER Q28(b)
\item
% >>> BEGIN ANSWER Q28(c)
% 在这里写英文推导或证明；完成后可删除下方空白占位命令。
\AnswerBlank
% <<< END ANSWER Q28(c)
\end{enumerate}
% END ANSWER Q28
\newpage\section*{Section D. Applications}
\textbf{29. Audit one binary batch update and the resulting decision}\hfill\textbf{[12 marks]}\par\medskip
\begin{enumerate}[label=(\alph*),leftmargin=8mm,itemsep=3pt]
\item Train on \((x,r)=(-2,0),(1,1),(3,1)\), starting from slope \(w=0\) and intercept \(b=0\). Use summed cross-entropy and one batch gradient step with learning rate \(0.1\). Compute the gradients and new parameters.\hfill\textbf{[4 marks]}
\item Compute the new probabilities and total cross-entropy. Compare it to the initial loss; exact logarithmic expressions are acceptable.\hfill\textbf{[4 marks]}
\item Using the equal-cost threshold, classify all three training points. Does a zero classification error here imply that further loss reduction is impossible?\hfill\textbf{[2 marks]}
\item Now suppose predicting \(C_1\) when the truth is \(C_2\) costs four, while the reverse error costs one, with zero correct-decision loss. Derive the posterior threshold and classify the same points under this cost rule.\hfill\textbf{[2 marks]}
\end{enumerate}\par\medskip\hrule\medskip\textbf{Student response}\par
% BEGIN ANSWER Q29 / 第29题答题区
\begin{enumerate}[label=(\alph*),leftmargin=8mm,itemsep=3pt]
\item
% >>> BEGIN ANSWER Q29(a)
% 在这里写英文推导或证明；完成后可删除下方空白占位命令。
\AnswerBlank
% <<< END ANSWER Q29(a)
\item
% >>> BEGIN ANSWER Q29(b)
% 在这里写英文推导或证明；完成后可删除下方空白占位命令。
\AnswerBlank
% <<< END ANSWER Q29(b)
\item
% >>> BEGIN ANSWER Q29(c)
% 在这里写英文推导或证明；完成后可删除下方空白占位命令。
\AnswerBlank
% <<< END ANSWER Q29(c)
\item
% >>> BEGIN ANSWER Q29(d)
% 在这里写英文推导或证明；完成后可删除下方空白占位命令。
\AnswerBlank
% <<< END ANSWER Q29(d)
\end{enumerate}
% END ANSWER Q29
\newpage\section*{Section D. Applications}
\textbf{30. Train a three-class model and check the normalization}\hfill\textbf{[12 marks]}\par\medskip
\begin{enumerate}[label=(\alph*),leftmargin=8mm,itemsep=3pt]
\item Consider one training sample \(x=(2,-1)^{\mathsf T}\) in class \(C_3\). Use augmented input \(\widetilde x=(1,2,-1)^{\mathsf T}\) and a row-weight matrix \(B\in\mathbb R^{3\times3}\), initially zero. Compute the one-hot target, probabilities, and loss.\hfill\textbf{[3 marks]}
\item Compute the full gradient matrix and update \(B\) using learning rate \(1/2\). Display every entry and state its dimension.\hfill\textbf{[5 marks]}
\item Recompute the logits, the probability of \(C_3\), and the new loss at this sample. Identify the predicted class.\hfill\textbf{[2 marks]}
\item Add the row vector \((10,0,0)\) to each fitted row of \(B\). Prove whether probabilities change, and describe a numerically stable way to compute the original softmax.\hfill\textbf{[2 marks]}
\end{enumerate}\par\medskip\hrule\medskip\textbf{Student response}\par
% BEGIN ANSWER Q30 / 第30题答题区
\begin{enumerate}[label=(\alph*),leftmargin=8mm,itemsep=3pt]
\item
% >>> BEGIN ANSWER Q30(a)
% 在这里写英文推导或证明；完成后可删除下方空白占位命令。
\AnswerBlank
% <<< END ANSWER Q30(a)
\item
% >>> BEGIN ANSWER Q30(b)
% 在这里写英文推导或证明；完成后可删除下方空白占位命令。
\AnswerBlank
% <<< END ANSWER Q30(b)
\item
% >>> BEGIN ANSWER Q30(c)
% 在这里写英文推导或证明；完成后可删除下方空白占位命令。
\AnswerBlank
% <<< END ANSWER Q30(c)
\item
% >>> BEGIN ANSWER Q30(d)
% 在这里写英文推导或证明；完成后可删除下方空白占位命令。
\AnswerBlank
% <<< END ANSWER Q30(d)
\end{enumerate}
% END ANSWER Q30
\end{document}
